\documentclass[10pt,a4paper]{amsart}
\usepackage{amsmath,amssymb,amsthm,graphicx}
% Retained minimal preamble for the Stacks Project chapter sites-cohomology.tex (Cohomology on Sites).
% Verbatim excerpt of the chapter preamble: Xy-pic setup (native preamble lines 14-18) and the native
% theorem declarations with the mathematical macros (native preamble lines 154-217). Copied mechanically;
% nothing added, removed or redefined. The wrapper class and the amsmath/amssymb/amsthm/graphicx base
% packages are supplied by the assembler, so the class/externaldocument part of the native preamble is omitted.
% For commutative diagrams we use Xy-pic
\usepackage[all]{xy}

% We use 2cell for 2-commutative diagrams.
\xyoption{2cell}
% Theorem environments.
%
\theoremstyle{plain}
\newtheorem{theorem}[subsection]{Theorem}
\newtheorem{proposition}[subsection]{Proposition}
\newtheorem{lemma}[subsection]{Lemma}

\theoremstyle{definition}
\newtheorem{definition}[subsection]{Definition}
\newtheorem{example}[subsection]{Example}
\newtheorem{exercise}[subsection]{Exercise}
\newtheorem{situation}[subsection]{Situation}

\theoremstyle{remark}
\newtheorem{remark}[subsection]{Remark}
\newtheorem{remarks}[subsection]{Remarks}

\numberwithin{equation}{subsection}

% Macros
%
\def\lim{\mathop{\mathrm{lim}}\nolimits}
\def\colim{\mathop{\mathrm{colim}}\nolimits}
\def\Spec{\mathop{\mathrm{Spec}}}
\def\Hom{\mathop{\mathrm{Hom}}\nolimits}
\def\Ext{\mathop{\mathrm{Ext}}\nolimits}
\def\SheafHom{\mathop{\mathcal{H}\!\mathit{om}}\nolimits}
\def\SheafExt{\mathop{\mathcal{E}\!\mathit{xt}}\nolimits}
\def\Sch{\mathit{Sch}}
\def\Mor{\mathop{\mathrm{Mor}}\nolimits}
\def\Ob{\mathop{\mathrm{Ob}}\nolimits}
\def\Sh{\mathop{\mathit{Sh}}\nolimits}
\def\NL{\mathop{N\!L}\nolimits}
\def\CH{\mathop{\mathrm{CH}}\nolimits}
\def\proetale{{pro\text{-}\acute{e}tale}}
\def\etale{{\acute{e}tale}}
\def\QCoh{\mathit{QCoh}}
\def\Ker{\mathop{\mathrm{Ker}}}
\def\Im{\mathop{\mathrm{Im}}}
\def\Coker{\mathop{\mathrm{Coker}}}
\def\Coim{\mathop{\mathrm{Coim}}}

% Boxtimes
%
\DeclareMathSymbol{\boxtimes}{\mathbin}{AMSa}{"02}

%
% Macros for moduli stacks/spaces
%
\def\QCohstack{\mathcal{QC}\!\mathit{oh}}
\def\Cohstack{\mathcal{C}\!\mathit{oh}}
\def\Spacesstack{\mathcal{S}\!\mathit{paces}}
\def\Quotfunctor{\mathrm{Quot}}
\def\Hilbfunctor{\mathrm{Hilb}}
\def\Curvesstack{\mathcal{C}\!\mathit{urves}}
\def\Polarizedstack{\mathcal{P}\!\mathit{olarized}}
\def\Complexesstack{\mathcal{C}\!\mathit{omplexes}}
% \Pic is the operator that assigns to X its picard group, usage \Pic(X)
% \Picardstack_{X/B} denotes the Picard stack of X over B
% \Picardfunctor_{X/B} denotes the Picard functor of X over B
\def\Pic{\mathop{\mathrm{Pic}}\nolimits}
\def\Picardstack{\mathcal{P}\!\mathit{ic}}
\def\Picardfunctor{\mathrm{Pic}}
\def\Deformationcategory{\mathcal{D}\!\mathit{ef}}
\title{The integral Cech chain complex of a covering family in a category is exact in positive degrees}
\author{The Stacks Project Authors}
\date{stacks-project:sites-cohomology.tex:local-snapshot-20261011; GNU FDL 1.2 or later}
\begin{document}
\maketitle
\noindent\textit{Editorial note.} All delivered mathematical text is the author's own, copied from the pinned Stacks Project chapter file sites-cohomology.tex; no statement or proof step was rewritten, reordered or removed, and the proof is complete in the source, with no omitted verification.\par\medskip
\noindent\textit{Editorial note.} The leading context span is the author's own Lemma lemma-cech-map-into (tag 03AS) of the same section, printed because the counted statement refers to the chain complex Z\_\{U,bullet\} defined there, including its signs (-1)\textasciicircum{}j and the placement of the last nonzero term in degree 0. That span defines the label lemma-cech-map-into, which is therefore not listed among the external references; its own proof is not printed because the counted proof does not use it.\par\medskip
\noindent\textit{Editorial note.} Because the counted statement and its printed context are self-contained with respect to the chapter, the card declares no external references at all; the two items of the standard prerequisite list are background notions, not cited labels.\par\medskip
\noindent\textit{Editorial note.} The preamble is the chapter's own retained minimal preamble (sources/stacks-sites-cohomology/retained-preamble.tex): a verbatim excerpt of the native Xy-pic setup and of the native theorem declarations and mathematical macros. The native class line, the xr-hyper externaldocument block and the reference/slogan/history display environments are omitted because the assembler supplies a plain amsart wrapper and the excerpt uses none of those environments; no macro was invented or redefined.\par\medskip
\noindent\textit{Allowed standard prerequisites (editorial).}\begin{itemize}
\item Free abelian presheaves represented by objects of the category, and the adjunction Hom(Z\_X, F) = F(X). Modules on Sites, Stacks tag 03A8 (sites-modules-lemma-obvious-adjointness), used by the printed context lemma for the identification recorded at its proof line 1021. Fixes the meaning of the terms Z\_\{U\_\{i0\} times\_U ... times\_U U\_\{ip\}\} out of which the chain complex of the statement is built; the exactness argument itself only uses free abelian groups on sets.
\item Abelian presheaves on a category, morphism sets and fibre products in an arbitrary category. the standing conventions of the chapter, restated inside the statement at source lines 1026-1029. The statement only assumes that the fibre products U\_\{i0\} times\_U ... times\_U U\_\{ip\} exist in the category C; the proof works with the sets Mor\_C(V, U\_\{i0\} times\_U ... times\_U U\_\{ip\}) and never needs a site structure.
\end{itemize}\medskip
\begin{lemma}
\label{lemma-cech-map-into}
Let $\mathcal{C}$ be a category. Let $\mathcal{U} = \{U_i \to U\}_{i \in I}$
be a family of morphisms with fixed target such that all fibre products
$U_{i_0} \times_U \ldots \times_U U_{i_p}$ exist in $\mathcal{C}$.
Consider the chain complex $\mathbf{Z}_{\mathcal{U}, \bullet}$
of abelian presheaves
$$
\ldots
\to
\bigoplus_{i_0i_1i_2} \mathbf{Z}_{U_{i_0} \times_U U_{i_1} \times_U U_{i_2}}
\to
\bigoplus_{i_0i_1} \mathbf{Z}_{U_{i_0} \times_U U_{i_1}}
\to
\bigoplus_{i_0} \mathbf{Z}_{U_{i_0}}
\to 0 \to \ldots
$$
where the last nonzero term is placed in degree $0$
and where the map
$$
\mathbf{Z}_{U_{i_0} \times_U \ldots \times_u U_{i_{p + 1}}}
\longrightarrow
\mathbf{Z}_{U_{i_0} \times_U
\ldots \widehat{U_{i_j}} \ldots \times_U U_{i_{p + 1}}}
$$
is given by $(-1)^j$ times the canonical map.
Then there is an isomorphism
$$
\Hom_{\textit{PAb}(\mathcal{C})}(\mathbf{Z}_{\mathcal{U}, \bullet}, \mathcal{F})
=
\check{\mathcal{C}}^\bullet(\mathcal{U}, \mathcal{F})
$$
functorial in $\mathcal{F} \in \Ob(\textit{PAb}(\mathcal{C}))$.
\end{lemma}

\begin{lemma}
\label{lemma-homology-complex}
Let $\mathcal{C}$ be a category. Let
$\mathcal{U} = \{f_i : U_i \to U\}_{i \in I}$ be a family of morphisms
with fixed target such that all fibre products
$U_{i_0} \times_U \ldots \times_U U_{i_p}$ exist in $\mathcal{C}$.
The chain complex $\mathbf{Z}_{\mathcal{U}, \bullet}$ of presheaves
of Lemma \ref{lemma-cech-map-into} above is exact in positive
degrees, i.e., the homology presheaves
$H_i(\mathbf{Z}_{\mathcal{U}, \bullet})$ are zero for $i > 0$.
\end{lemma}

\begin{proof}
Let $V$ be an object of $\mathcal{C}$. We have to show that the chain complex
of abelian groups $\mathbf{Z}_{\mathcal{U}, \bullet}(V)$ is exact in
degrees $> 0$. This is the complex
$$
\xymatrix{
\ldots \ar[d] \\
\bigoplus_{i_0i_1i_2}
\mathbf{Z}[
\Mor_\mathcal{C}(V, U_{i_0} \times_U U_{i_1} \times_U U_{i_2})
]
\ar[d] \\
\bigoplus_{i_0i_1}
\mathbf{Z}[
\Mor_\mathcal{C}(V, U_{i_0} \times_U U_{i_1})
]
\ar[d] \\
\bigoplus_{i_0}
\mathbf{Z}[
\Mor_\mathcal{C}(V, U_{i_0})
] \ar[d] \\
0
}
$$
For any morphism $\varphi : V \to U$ denote
$\Mor_\varphi(V, U_i) = \{\varphi_i : V \to U_i \mid
f_i \circ \varphi_i = \varphi\}$. We will use a similar notation
for $\Mor_\varphi(V, U_{i_0} \times_U \ldots \times_U U_{i_p})$.
Note that composing with the various projection maps between the
fibred products $U_{i_0} \times_U \ldots \times_U U_{i_p}$ preserves
these morphism sets. Hence we see that the complex above
is the same as the complex
$$
\xymatrix{
\ldots \ar[d] \\
\bigoplus_\varphi
\bigoplus_{i_0i_1i_2}
\mathbf{Z}[
\Mor_\varphi(V, U_{i_0} \times_U U_{i_1} \times_U U_{i_2})
]
\ar[d] \\
\bigoplus_\varphi
\bigoplus_{i_0i_1}
\mathbf{Z}[
\Mor_\varphi(V, U_{i_0} \times_U U_{i_1})
]
\ar[d] \\
\bigoplus_\varphi
\bigoplus_{i_0}
\mathbf{Z}[
\Mor_\varphi(V, U_{i_0})
] \ar[d] \\
0
}
$$
Next, we make the remark that we have
$$
\Mor_\varphi(V, U_{i_0} \times_U \ldots \times_U U_{i_p})
=
\Mor_\varphi(V, U_{i_0}) \times \ldots
\times \Mor_\varphi(V, U_{i_p})
$$
Using this and the fact that $\mathbf{Z}[A] \oplus \mathbf{Z}[B] =
\mathbf{Z}[A \amalg B]$ we see that the complex becomes
$$
\xymatrix{
\ldots \ar[d] \\
\bigoplus_\varphi
\mathbf{Z}\left[
\coprod_{i_0i_1i_2}
\Mor_\varphi(V, U_{i_0}) \times \Mor_\varphi(V, U_{i_1}) \times
\Mor_\varphi(V, U_{i_2})
\right]
\ar[d] \\
\bigoplus_\varphi
\mathbf{Z}\left[
\coprod_{i_0i_1}
\Mor_\varphi(V, U_{i_0}) \times \Mor_\varphi(V, U_{i_1})
\right]
\ar[d] \\
\bigoplus_\varphi
\mathbf{Z}\left[
\coprod_{i_0}
\Mor_\varphi(V, U_{i_0})
\right] \ar[d] \\
0
}
$$
Finally, on setting $S_\varphi = \coprod_{i \in I} \Mor_\varphi(V, U_i)$
we see that we get
$$
\bigoplus\nolimits_\varphi \left(\ldots \to
\mathbf{Z}[S_\varphi \times S_\varphi \times S_\varphi] \to
\mathbf{Z}[S_\varphi \times S_\varphi] \to
\mathbf{Z}[S_\varphi] \to 0 \to \ldots
\right)
$$
Thus we have simplified our task. Namely, it suffices to show that
for any nonempty set $S$ the (extended) complex of free abelian groups
$$
\ldots \to
\mathbf{Z}[S \times S \times S] \to
\mathbf{Z}[S \times S] \to
\mathbf{Z}[S] \xrightarrow{\Sigma} \mathbf{Z} \to 0 \to \ldots
$$
is exact in all degrees. To see this fix an element $s \in S$, and
use the homotopy
$$
n_{(s_0, \ldots, s_p)} \longmapsto n_{(s, s_0, \ldots, s_p)}
$$
with obvious notations.
\end{proof}

\end{document}
